\documentclass[11pt,a4paper]{article}
\usepackage[utf8]{inputenc}
\usepackage[T1]{fontenc}
\usepackage[spanish,es-nodecimaldot]{babel}
\usepackage{amsmath,amssymb}
\usepackage{booktabs}
\usepackage{tikz}
\usetikzlibrary{babel}
\usepackage[margin=2.3cm]{geometry}
\usepackage{float}
\usepackage{siunitx}
\sisetup{output-decimal-marker={,}}

\newcommand{\kNmm}{\si{kN/mm}}
\newcommand{\kNmmm}{\si{kN\cdot mm}}
\newcommand{\Eeq}{\si{kN/mm^2}}

\usepackage{fancyhdr}
\pagestyle{fancy}\fancyhf{}
\lhead{Materiales Aeroespaciales}\rhead{Fecha: 15/10/2026}\cfoot{\thepage}
\usepackage{titlesec}
\titleformat{\section}{\normalfont\Large\bfseries}{}{0pt}{}
\titleformat{\subsection}{\normalfont\large\bfseries}{}{0pt}{}
\titleformat{\subsubsection}{\normalfont\normalsize\bfseries}{}{0pt}{}
\setlength{\parskip}{6pt}\setlength{\parindent}{0pt}

\begin{document}
\begin{titlepage}
\centering
\vspace*{3cm}
{\Huge Materiales Aeroespaciales\par}
\vspace{0.6cm}
{\LARGE TP: Laminados especiales\par}
\vspace{2cm}
{\large Profesores: Ing. Anibal Cesar Vettorel\\
\hspace{2.2cm}Ing. Carlos Mario Carbonel\par}
\vspace{1cm}
{\large Alumno: Gimenez, Caetano -- 27446\par}
\vspace{1cm}
{\large Fecha: 15/10/2026\par}
{\large Año: 2026\par}
\end{titlepage}


\section{Introducción;}

El presente trabajo tiene como objetivo aplicar los conceptos de la Teoría Clásica de Laminados al análisis de distintas configuraciones de laminados de carbono/epoxi. Para cada una se calcularon las matrices $A$, $B$ y $D$ y las constantes elásticas equivalentes, tanto para el comportamiento en membrana como en flexión, y se analizaron los resultados a partir de la simetría, la orientación y la secuencia de apilamiento.

En la primera parte se estudió un laminado cuasi-isotrópico $(0/45/-45/90)_s$. En la segunda, tres laminados cross-ply: $(0/0/90/90)$, $(0/90/0/90)$ y $(0/90)_s$. Los dos primeros son antisimétricos y el tercero es simétrico, que se incluyó a modo de comparación. Todos los cálculos se hicieron con un script de Python que sigue las mismas ecuaciones que se detallan más abajo.

\section{Datos e hipótesis;}

La lámina es de carbono/epoxi de alta resistencia, con un espesor $t=0{,}125\ \mathrm{mm}$. Las constantes elásticas son
\[
E_1=140,\quad E_2=10,\quad G_{12}=5\ \Eeq,\qquad \nu_{12}=0{,}3,
\qquad \nu_{21}=\nu_{12}\frac{E_2}{E_1}=0{,}0214 .
\]
Se trabajó en unidades de kN y mm. El sistema de coordenadas es el del enunciado: $x$ e $y$ en el plano del laminado y $z$ positivo hacia arriba, con el origen en el plano medio. Los ángulos se miden desde $x$ en sentido antihorario y la lámina 1 es la que está más abajo ($z$ más negativo).

Se adoptaron las hipótesis habituales de la teoría: las láminas están perfectamente adheridas entre sí, cada una trabaja en tensión plana ($\sigma_z=\tau_{xz}=\tau_{yz}=0$) con comportamiento elástico lineal, y las normales al plano medio se mantienen rectas y perpendiculares tras la deformación (Kirchhoff). Esto implica una deformación lineal en el espesor, $\varepsilon=\varepsilon^0+z\kappa$.

Las rigideces reducidas de la lámina en sus ejes son
\[
Q_{11}=\frac{E_1}{1-\nu_{12}\nu_{21}},\quad
Q_{22}=\frac{E_2}{1-\nu_{12}\nu_{21}},\quad
Q_{12}=\frac{\nu_{12}E_2}{1-\nu_{12}\nu_{21}},\quad
Q_{66}=G_{12},
\]
y para una lámina a un ángulo $\theta$ se obtiene $[\bar Q]=[T]^{-1}[Q][R][T][R]^{-1}$, con $c=\cos\theta$, $s=\sin\theta$, $[R]=\mathrm{diag}(1,1,2)$ y
\[
[T]=\begin{bmatrix}c^2&s^2&2cs\\ s^2&c^2&-2cs\\ -cs&cs&c^2-s^2\end{bmatrix}.
\]
Las matrices del laminado se obtienen sumando lámina a lámina, siendo $z_k$ la cota de la cara superior de la lámina $k$:
\[
A_{ij}=\sum_k \bar Q_{ij}^{(k)}(z_k-z_{k-1}),\quad
B_{ij}=\frac12\sum_k \bar Q_{ij}^{(k)}(z_k^2-z_{k-1}^2),\quad
D_{ij}=\frac13\sum_k \bar Q_{ij}^{(k)}(z_k^3-z_{k-1}^3),
\]
\[
\begin{Bmatrix}N\\M\end{Bmatrix}=
\begin{bmatrix}A&B\\B&D\end{bmatrix}
\begin{Bmatrix}\varepsilon^0\\\kappa\end{Bmatrix}.
\]
Las constantes equivalentes se calcularon con $a=A^{-1}$ para membrana y $d=D^{-1}$ para flexión, siendo $h$ el espesor total del laminado:
\begin{align}
\text{Membrana:}&\quad
E_x=\frac{1}{h\,a_{11}},\ E_y=\frac{1}{h\,a_{22}},\ G_{xy}=\frac{1}{h\,a_{66}},\
\nu_{xy}=-\frac{a_{12}}{a_{11}},\\
\text{Flexión:}&\quad
E^f_x=\frac{12}{h^3 d_{11}},\ E^f_y=\frac{12}{h^3 d_{22}},\ G^f_{xy}=\frac{12}{h^3 d_{66}},\
\nu^f_{xy}=-\frac{d_{12}}{d_{11}}.
\end{align}

Con los datos de la lámina resultan las siguientes rigideces, en $\Eeq$:
\[
[Q]=[\bar Q]_{0^\circ}=\begin{bmatrix}140{,}906&3{,}019&0\\3{,}019&10{,}065&0\\0&0&5\end{bmatrix},\qquad
[\bar Q]_{90^\circ}=\begin{bmatrix}10{,}065&3{,}019&0\\3{,}019&140{,}906&0\\0&0&5\end{bmatrix},
\]
\[
[\bar Q]_{45^\circ}=\begin{bmatrix}44{,}252&34{,}252&32{,}710\\34{,}252&44{,}252&32{,}710\\32{,}710&32{,}710&36{,}233\end{bmatrix},\qquad
[\bar Q]_{-45^\circ}=\begin{bmatrix}44{,}252&34{,}252&-32{,}710\\34{,}252&44{,}252&-32{,}710\\-32{,}710&-32{,}710&36{,}233\end{bmatrix}.
\]

\section{Parte 1: laminado cuasi-isotrópico;}

\subsection{Definición y condiciones;}
Se entiende por laminado cuasi-isotrópico a aquel cuya matriz $[A]$ es isótropa, es decir $A_{11}=A_{22}$, $A_{16}=A_{26}=0$ y $A_{66}=(A_{11}-A_{12})/2$. De esta forma las propiedades en el plano no dependen de la dirección. Para lograrlo la configuración debe tener tres o más orientaciones, con los ángulos igualmente espaciados en $\pi/N$ (por ejemplo $0/60/-60$ o $0/45/90/-45$) y con el mismo espesor total en cada una. Hay que aclarar que esta isotropía vale solo en membrana, ya que en flexión la matriz $[D]$ depende de la secuencia de apilamiento y en general no es isótropa.

\subsection{Esquema;}
El laminado $(0/45/-45/90)_s$ tiene 8 láminas y un espesor total de $h=1{,}0$ mm. Es simétrico respecto del plano medio, por lo que las láminas 5 a 8 repiten en espejo a las láminas 4 a 1.

\begin{figure}[H]
\centering
\begin{tikzpicture}[scale=0.9]
  \foreach \i/\ang in {1/0,2/45,3/{-45},4/90,5/90,6/{-45},7/45,8/0}{
    \pgfmathsetmacro{\y}{(\i-1)*0.55}
    \draw[fill=gray!15] (0,\y) rectangle (4,\y+0.55);
    \node at (2,\y+0.275) {Lámina \i: $\ang^\circ$};
  }
  \draw[dashed] (-0.4,2.2) -- (4.4,2.2) node[right] {plano medio ($z=0$)};
  \draw[->,thick] (-0.9,0) -- (-0.9,4.9) node[above] {$z$};
  \node[left] at (-0.9,0.28) {$z=-0{,}5$};
  \node[left] at (-0.9,4.4) {$z=+0{,}5$};
\end{tikzpicture}
\caption{Laminado $(0/45/-45/90)_s$. La lámina 1 está abajo y $z$ es positivo hacia arriba (cotas en mm).}
\end{figure}

\subsection{Matrices A, B y D;}
\[
[A]=\begin{bmatrix}59{,}869&18{,}636&0\\18{,}636&59{,}869&0\\0&0&20{,}616\end{bmatrix}\kNmm,\qquad
[B]=[0]\ \si{kN},
\]
\[
[D]=\begin{bmatrix}8{,}2997&1{,}3090&0{,}5111\\1{,}3090&2{,}1665&0{,}5111\\0{,}5111&0{,}5111&1{,}4740\end{bmatrix}\kNmmm .
\]

Como se puede ver, $[A]$ cumple las condiciones de isotropía: $A_{11}=A_{22}$, $A_{16}=A_{26}=0$ y $A_{66}=(59{,}869-18{,}636)/2=20{,}616$, y por lo tanto no hay acoplamiento entre tracción y corte. Esto es consistente con lo esperado para un cuasi-isotrópico. La matriz $[B]$ es nula porque el laminado es simétrico respecto del plano medio, por lo que una carga en el plano no genera curvatura.

Por otro lado, $[D]$ no es isótropa. Ocurre que $D_{11}\neq D_{22}$ debido a que las láminas a $0^\circ$ están en las caras exteriores, donde tienen mayor brazo respecto del plano medio, mientras que las de $90^\circ$ están más cerca del centro. Además, $D_{16}$ y $D_{26}$ son distintos de cero por la presencia de las láminas a $\pm45^\circ$, lo que genera acoplamiento entre flexión y torsión. Es decir que el laminado es cuasi-isotrópico en membrana pero no en flexión.

\subsection{Constantes elásticas equivalentes;}
\begin{center}
\begin{tabular}{lcccc}
\toprule
 & $E_x\ [\Eeq]$ & $E_y\ [\Eeq]$ & $G_{xy}\ [\Eeq]$ & $\nu_{xy}$ \\
\midrule
Membrana & 54,07 & 54,07 & 20,62 & 0,311 \\
Flexión  & 89,74 & 21,98 & 16,18 & 0,569 \\
\bottomrule
\end{tabular}
\end{center}
En membrana se verifica la relación de un material isótropo, $G=E/[2(1+\nu)]=54{,}07/2{,}622=20{,}62\ \Eeq$. En flexión, en cambio, las constantes son muy distintas entre $x$ e $y$ (con $\nu^f_{yx}=0{,}139$), lo que muestra que a flexión el laminado se comporta como ortótropo y es sensible a la secuencia.

\section{Parte 2: laminados cross-ply antisimétricos;}

\subsection{Definición;}
Un laminado cross-ply es aquel que solo tiene láminas a $0^\circ$ y $90^\circ$. Se lo llama antisimétrico cuando, para cada lámina de orientación $\theta$ a una distancia $z$ del plano medio, existe otra de orientación $90^\circ-\theta$ a una distancia $-z$, con el mismo material y espesor. En otras palabras, el apilado por encima del plano medio es el complemento del de abajo. La consecuencia es que $A_{16}=A_{26}=D_{16}=D_{26}=0$ y que $B_{11}=-B_{22}\neq0$, con el resto de $[B]$ nulo. Esto genera acoplamiento entre membrana y flexión: una tracción en $x$ curva el laminado.

Los laminados A $(0/0/90/90)$ y B $(0/90/0/90)$ son antisimétricos. El C $(0/90)_s=[0/90/90/0]$ es simétrico y lo incluimos como referencia.

\subsection{Esquemas;}
\begin{figure}[H]
\centering
\begin{tikzpicture}[scale=0.8]
  \foreach \x/\lab/\a/\b/\c/\d in {0/A/0/0/90/90, 6/B/0/90/0/90, 12/C/0/90/90/0}{
    \foreach \i/\ang in {1/\a,2/\b,3/\c,4/\d}{
      \pgfmathsetmacro{\y}{(\i-1)*0.6}
      \draw[fill=gray!15] (\x,\y) rectangle (\x+3.2,\y+0.6);
      \node at (\x+1.6,\y+0.3) {\i: $\ang^\circ$};
    }
    \draw[dashed] (\x-0.2,1.2) -- (\x+3.4,1.2);
    \node at (\x+1.6,-0.5) {Laminado \lab};
  }
  \draw[->,thick] (-0.7,0) -- (-0.7,2.8) node[above] {$z$};
\end{tikzpicture}
\caption{Configuraciones A $(0/0/90/90)$, B $(0/90/0/90)$ y C $(0/90)_s$. Cada lámina es de $0{,}125$ mm, $h=0{,}5$ mm, $z$ es positivo hacia arriba y la lámina 1 está abajo (cotas $z=-0{,}25;-0{,}125;0;0{,}125;0{,}25$ mm).}
\end{figure}

\subsection{Matrices A, B y D;}
Las tres configuraciones tienen la misma cantidad de láminas a $0^\circ$ y a $90^\circ$, por lo que comparten la matriz
\[
[A]=\begin{bmatrix}37{,}743&1{,}510&0\\1{,}510&37{,}743&0\\0&0&2{,}500\end{bmatrix}\kNmm .
\]
Lo que cambia entre ellas son $[B]$ (en \si{kN}) y $[D]$ (en \kNmmm):
\[
\text{A: }[B]=\begin{bmatrix}-4{,}0888&0&0\\0&4{,}0888&0\\0&0&0\end{bmatrix},\quad
[D]=\begin{bmatrix}0{,}7863&0{,}0315&0\\0{,}0315&0{,}7863&0\\0&0&0{,}0521\end{bmatrix}
\]
\[
\text{B: }[B]=\begin{bmatrix}-2{,}0444&0&0\\0&2{,}0444&0\\0&0&0\end{bmatrix},\quad
[D]=\begin{bmatrix}0{,}7863&0{,}0315&0\\0{,}0315&0{,}7863&0\\0&0&0{,}0521\end{bmatrix}
\]
\[
\text{C: }[B]=[0],\quad
[D]=\begin{bmatrix}1{,}2974&0{,}0315&0\\0{,}0315&0{,}2752&0\\0&0&0{,}0521\end{bmatrix}
\]

En los tres casos $A_{11}=A_{22}$ y $A_{16}=A_{26}=0$, es decir que son laminados balanceados, ortótropos en membrana y sin acoplamiento entre tracción y corte. En A y B aparece $B_{11}=-B_{22}\neq0$, que es el acoplamiento entre membrana y flexión que se esperaba por ser antisimétricos. El valor de $B_{11}$ en A es el doble que en B: en A las láminas a $0^\circ$ están agrupadas en una misma cara y el efecto se acumula, mientras que en B el apilado alternado lo compensa en parte.

El laminado C es simétrico, por lo que $[B]=0$ y no hay acoplamiento. Sin embargo, su matriz $[D]$ está muy desbalanceada ($D_{11}\gg D_{22}$) porque las láminas a $0^\circ$ están en las caras exteriores y las de $90^\circ$ en el interior. Por último, $[D]$ resulta igual en A y en B, ya que en ambos casos las láminas de cada orientación quedan en cotas simétricas respecto de las de la otra orientación.

\subsection{Constantes elásticas equivalentes;}
Primero se calcularon con $A^{-1}$ y $D^{-1}$, sin considerar el acoplamiento:
\begin{center}
\begin{tabular}{llcccc}
\toprule
Laminado & Modo & $E_x\ [\Eeq]$ & $E_y\ [\Eeq]$ & $G_{xy}\ [\Eeq]$ & $\nu_{xy}$ \\
\midrule
A & membrana & 75,36 & 75,36 & 5,00 & 0,040 \\
A & flexión  & 75,36 & 75,36 & 5,00 & 0,040 \\
B & membrana & 75,36 & 75,36 & 5,00 & 0,040 \\
B & flexión  & 75,36 & 75,36 & 5,00 & 0,040 \\
C & membrana & 75,36 & 75,36 & 5,00 & 0,040 \\
C & flexión  & 124,21 & 26,35 & 5,00 & 0,114 \\
\bottomrule
\end{tabular}
\end{center}

En A y B, como $[B]\neq0$, estos valores sobreestiman la rigidez real. Por eso se invirtió también la matriz completa $\begin{bmatrix}A&B\\B&D\end{bmatrix}$ y se calcularon las constantes con los bloques $a_{ij}$ y $d_{ij}$ de la flexibilidad resultante:
\begin{center}
\begin{tabular}{lcccc}
\toprule
Laminado & $E_x=E_y\ [\Eeq]$ (membrana y flexión) & $G_{xy}\ [\Eeq]$ & $\nu_{xy}$ & Reducción de $E$ \\
\midrule
A & 32,84 & 5,00 & 0,040 & 56\,\% \\
B & 64,73 & 5,00 & 0,040 & 14\,\% \\
\bottomrule
\end{tabular}
\end{center}
La pérdida de rigidez es mucho mayor en A que en B, lo que va de la mano con el valor de $B_{11}$ de cada uno.

\section{Comparación y conclusiones;}

\begin{center}\small
\setlength{\tabcolsep}{3pt}
\begin{tabular}{lcccc}
\toprule
Laminado & Simetría & $B$ & $A$ & $D$ \\
\midrule
$(0/45/-45/90)_s$ & simétrico & $0$ & isótropa & ortótropa con $D_{16},D_{26}\neq0$ \\
A $(0/0/90/90)$ & antisimétrico & $B_{11}=-B_{22}=-4{,}09$ & ortótropa & ortótropa, $D_{11}{=}D_{22}$ \\
B $(0/90/0/90)$ & antisimétrico & $B_{11}=-B_{22}=-2{,}04$ & ortótropa & ortótropa, $D_{11}{=}D_{22}$ \\
C $(0/90)_s$ & simétrico & $0$ & ortótropa & ortótropa, $D_{11}\gg D_{22}$ \\
\bottomrule
\end{tabular}
\end{center}

A modo de conclusión, el laminado cuasi-isotrópico logra propiedades iguales en todas las direcciones del plano ($E=54{,}07$ y $G=20{,}62\ \Eeq$), pero a costa de tener menos rigidez en la dirección de las fibras que los cross-ply ($E_x=75{,}36\ \Eeq$), y además conserva acoplamiento entre flexión y torsión por las láminas a $\pm45^\circ$.

En cuanto a la simetría, un laminado simétrico respecto del plano medio tiene $[B]=0$ y no presenta acoplamiento entre membrana y flexión. En cambio, la antisimetría de A y B introduce ese acoplamiento y reduce la rigidez efectiva, 56\,\% en A y 14\,\% en B. En la práctica esto se traduce en alabeo de la pieza después del curado y bajo carga.

Por último, las tres configuraciones cross-ply tienen la misma rigidez de membrana, pero la secuencia de apilado define la rigidez a flexión: en C los $0^\circ$ están afuera y su rigidez a flexión en $x$ es 1,65 veces la de A y B, a cambio de una $E^f_y$ mucho menor. De esto se desprende que, para evitar el alabeo, conviene usar laminados simétricos, y que la cantidad y orientación de las láminas definen el comportamiento en membrana mientras que la secuencia define el de flexión.

\end{document}
